\section{Limitations and Capability Boundaries}
\label{sec:limitations}

% Honest accounting of what the system cannot do.

% 1. Existential vs universal detection
%   Existential detection (op(e,x,x) for SOME x) gives false positives
%   on non-group structures. Universal detection (double negation) is
%   correct but requires two strata of negation.

% 2. omega-rule defense
%   Pure-transfer rules (R(x)|-S(x)) are blocked at pattern level.
%   This is conservative: some correct patterns are lost (e.g., nonneg(_0)).

% 3. Backward application disabled
%   Apply rules (has_property -> formula) cause cascading pattern facts
%   in algebraic contexts. Currently forward-only detection.

% 4. Zorn's Lemma
%   System can verify Zorn on specific finite posets but cannot prove
%   the universal theorem. Requires second-order quantification and
%   structural induction, beyond first-order Horn + stratified negation.
%   The iterative chain construction produces depth-bounded "maximal"
%   elements that are artifacts of the depth limit, not genuine maximality.
%   Proof by contradiction correctly distinguishes real from fake maximality.

% 5. Property implication is inductive
%   implies_observed is based on current data (ext(P) subset ext(Q)).
%   Not a logical proof. Could be invalidated by new elements.
%   Separated from deductive implies to prevent modus ponens cascading.

% 6. Carrier size 3
%   All algebraic results are on 3-element carriers.
%   Transfer to Z4 verified (6/6), but systematic n>3 exploration
%   is computationally expensive.
